% arXiv submission draft.
% Compile: pdflatex paper.tex (run twice for cross-references).
% TODO before submission:
%   - Fill in \author{} and \thanks{} with real identity / affiliation / email.
%   - Update \date{} if needed.
%   - Verify all URLs.

\documentclass[11pt]{article}

\usepackage[utf8]{inputenc}
\usepackage[T1]{fontenc}
\usepackage[english]{babel}
\usepackage{amsmath, amssymb, amsthm}
\usepackage[margin=1in]{geometry}
\usepackage{hyperref}
\usepackage{booktabs}
\usepackage{array}
\usepackage{listings}
\usepackage{xcolor}
\usepackage{microtype}

\hypersetup{
  colorlinks=true,
  linkcolor=blue!50!black,
  citecolor=blue!50!black,
  urlcolor=blue!50!black,
}

\newtheorem{theorem}{Theorem}
\newtheorem{corollary}{Corollary}

\lstdefinestyle{pythonstyle}{
  language=Python,
  basicstyle=\ttfamily\footnotesize,
  keywordstyle=\color{blue!60!black},
  commentstyle=\color{green!40!black}\itshape,
  stringstyle=\color{red!60!black},
  showstringspaces=false,
  breaklines=true,
  frame=single,
  framerule=0.4pt,
  rulecolor=\color{black!30},
  numbers=none,
  xleftmargin=0.5em,
  xrightmargin=0.5em,
  aboveskip=0.5em,
  belowskip=0.5em,
}

\lstdefinestyle{pseudostyle}{
  basicstyle=\ttfamily\footnotesize,
  showstringspaces=false,
  breaklines=true,
  frame=single,
  framerule=0.4pt,
  rulecolor=\color{black!30},
  numbers=none,
  xleftmargin=0.5em,
  xrightmargin=0.5em,
  aboveskip=0.5em,
  belowskip=0.5em,
}

\title{A Noise-Invariant Determinism Theorem for Multi-Base \\
Post-Processing in Shor's Order Finding}

\author{%
  Author Name%
  \thanks{TODO: fill in author name, affiliation, email before submission. Code repository: \url{https://github.com/Hashevolution/shor}.}%
}

\date{June 2026}

\begin{document}
\maketitle

\begin{abstract}
We give an elementary observation about classical post-processing of multiple-base order-finding measurements in Shor's algorithm. Maintaining an accumulated exponent candidate $L$ --- the least common multiple of orders recovered from previous bases --- and augmenting standard continued-fraction post-processing with a divisor search over $L$ yields a procedure (which we call $(C)$) that recovers the order $r_a$ of any new base $a$ deterministically whenever $r_a \mid L$, \emph{independent of the measurement distribution}. As a corollary, the procedure is robust to any measurement-layer noise model: a measurement returning a value of $k$ drawn from any distribution still yields the correct order, provided the accumulated $L$ covers $r_a$. We verify this experimentally across 11 noise setups (depolarizing, readout flip, bias, phase decoherence, amplitude damping, modular exponentiation error) at three composite sizes and 500 trials each (16{,}500 measurements), and additionally at three larger semiprimes up to $N = 4087$ (1{,}200 further measurements) --- 17{,}700 total, zero violations. We discuss the theorem's relation to prior work --- Knill's lcm trick (1995), McAnally's larger-$Q$ convergent enumeration (2001), the Bach--Shallit textbook treatment of $r \mid \lambda(N)$, and the quantum algorithm for computing the Carmichael function (2021) --- and conclude that while every individual ingredient is folklore, the \emph{clean statement of the noise-invariance corollary} appears not to be made explicit in surveyed literature. We additionally show that two natural attempts to improve measurement count beyond this scheme --- adaptive base selection and lattice-based joint post-processing --- yield no further reduction, suggesting the theorem captures the essential structure of the problem at this scale.
\end{abstract}

\section{Background and Notation}
\label{sec:background}

Let $N$ be a composite integer, not a prime power, with $N$ odd. Define $Q = 2^t$ where $t = 2\lceil \log_2 N \rceil$ (so $Q \geq N^2$). For $a \in (\mathbb{Z}/N\mathbb{Z})^*$, let $r_a$ denote the multiplicative order of $a$ modulo $N$ (the smallest positive integer with $a^{r_a} \equiv 1 \bmod N$). Shor's order-finding~\cite{shor1994,nielsenchuang2000} measures $k \in \{0, 1, \ldots, Q-1\}$ from a distribution concentrated near $j \cdot Q / r_a$ for integer $j$. Standard classical post-processing uses the continued-fraction expansion of $k/Q$ to recover $r_a$.

Two classical facts are central:

\begin{description}
\item[(F1) Carmichael bound.] For all $a \in (\mathbb{Z}/N\mathbb{Z})^*$, $r_a \mid \lambda(N)$ where $\lambda$ is the Carmichael function --- the exponent of the group $(\mathbb{Z}/N\mathbb{Z})^*$.
\item[(F2) Order as minimum.] $r_a$ is by definition the smallest positive integer $d$ with $a^d \equiv 1 \bmod N$. Any $d$ satisfying this is a positive multiple of $r_a$.
\end{description}

\section{The $(C)$ Procedure}
\label{sec:procedure}

Let $L \in \mathbb{N}$ be a quantity maintained by the algorithm --- initially $L = 1$, updated as orders are recovered (see Algorithm~\ref{alg:1} below). For a measurement $k \in \{0, \ldots, Q-1\}$ and current $L$, define the candidate pool:
\[
  \mathrm{candidates}(k, L) \;:=\; \mathrm{convergents}(k/Q) \,\cup\, \mathrm{divisors}(L)
\]
where $\mathrm{convergents}(k/Q)$ is the set of denominators of continued-fraction convergents of $k/Q$ up to $N-1$, and $\mathrm{divisors}(L)$ is the set of positive divisors of $L$.

The $(C)$ post-processing is:

\begin{lstlisting}[style=pythonstyle]
def C(a, N, k, Q, L):
    cands = convergents(k/Q, max_denom=N-1) | divisors(L)
    valid = {d in cands : a^d == 1 mod N}
    if valid is empty: return 0
    return minimize_order(a, N, min(valid))
\end{lstlisting}

\noindent
where \texttt{minimize\_order(a, N, d)} reduces $d$ to its smallest positive divisor still satisfying $a^? \equiv 1$ (by removing small prime factors). See Appendix~\ref{app:impl} for the reference implementation.

\begin{lstlisting}[style=pseudostyle, caption={Multi-base accumulation.}, label={alg:1}]
L <- 1
for iteration in 1..max_iter:
    pick a in {2,...,N-1} uniformly at random
    if gcd(a, N) > 1: return factor gcd(a, N)
    if L > 1 and a^L == 1 mod N:
        # fast path: recover r_a classically
        r <- min{d divides L : a^d == 1 mod N}
    else:
        # slow path: quantum measurement followed by (C)
        k <- measure(a, N)
        r <- C(a, N, k, Q, L)
    if r == 0: continue
    L <- lcm(L, r)
    attempt factoring with (a, r) via standard Shor reduction
    or via factor_from_exponent(L) [Miller-Rabin style]
\end{lstlisting}

\section{Main Theorem}
\label{sec:theorem}

\begin{theorem}[$(C)$-determinism]\label{thm:c-det}
If $r_a \mid L$, then $C(a, N, k, Q, L) = r_a$ for every $k \in \{0, \ldots, Q-1\}$.
\end{theorem}

\begin{proof}
Assume $r_a \mid L$. Then $r_a \in \mathrm{divisors}(L) \subseteq \mathrm{candidates}(k, L)$. By construction, $a^{r_a} \equiv 1 \bmod N$, so $r_a \in \mathrm{valid}$. Hence $\mathrm{valid} \neq \emptyset$. By (F2), every $d \in \mathrm{valid}$ is a positive multiple of $r_a$, so $r_a \leq d$. Combining with $r_a \in \mathrm{valid}$ gives $\min(\mathrm{valid}) = r_a$. Finally, $\texttt{minimize\_order}(a, N, r_a) = r_a$ since $r_a$ is already the order.
\end{proof}

\begin{corollary}[Noise invariance]\label{cor:noise}
Theorem~\ref{thm:c-det} makes no assumption on the distribution from which $k$ is drawn. Hence: under any noise model that affects only the measurement step --- including arbitrary distortions of the QFT output distribution --- $(C)$ returns $r_a$ whenever $r_a \mid L$.
\end{corollary}

\begin{corollary}[$L$ integrity]\label{cor:integrity}
When Algorithm~\ref{alg:1} updates $L \leftarrow \mathrm{lcm}(L, r)$, the value $r$ returned by $(C)$ is always either $r_a$ (success) or $0$ (failure). In the success case, $r_a \mid \lambda(N)$ by (F1), so $L$ always remains a divisor of $\lambda(N)$. Under noise, $L$ cannot be corrupted to a value outside the divisors of $\lambda(N)$.
\end{corollary}

\begin{corollary}[Failure region]\label{cor:failure}
$(C)$ can fail only when $r_a \nmid L$. In that case, success depends on $\mathrm{convergents}(k/Q)$ containing some positive multiple of $r_a$ --- the classical ``$(B)$-path.''
\end{corollary}

\section{Empirical Verification}
\label{sec:empirical}

We verify Theorem~\ref{thm:c-det} across diverse noise models. For each measurement, we log $(r_a, L_{\mathrm{before}}, \mathrm{condition}, \mathrm{success})$ where $\mathrm{condition} := (r_a \mid L_{\mathrm{before}})$ and $\mathrm{success} := ((C) \text{ output} = r_a)$. We count:
\begin{itemize}
  \item \texttt{violations} := count of $\mathrm{condition} = \text{True} \land \mathrm{success} = \text{False}$ (theorem predicts $0$)
  \item \texttt{lucky} := $\mathrm{condition} = \text{False} \land \mathrm{success} = \text{True}$ (recovery via convergents alone)
  \item \texttt{covered} := $\mathrm{condition} = \text{True}$
\end{itemize}

Noise models tested (each implemented in \texttt{noise.py} of the companion code):

\begin{center}
\begin{tabular}{lll}
\toprule
Category & Model & Parameter \\
\midrule
A (measurement only) & depolarizing & $p \in [0, 1]$ \\
A & readout bit flip & $p$ per bit \\
A & bias-to-zero (adversarial) & $p$ \\
B (QFT input) & phase decoherence & $\sigma$ (rad) on amplitude \\
B & amplitude damping (T1) & $\gamma$ decay rate \\
C (structural) & modular exponentiation error & $q$ per $f(x)$ value \\
\bottomrule
\end{tabular}
\end{center}

\paragraph{Result.}
Across $N \in \{77, 143, 209\}$, 11 noise setups, and 500 trials each ($16{,}500$ measurements total): \textbf{violations $= 0$}. The theorem holds in every case.

The \texttt{lucky} count is small ($1$--$4$ per setup), confirming most successes flow through the \texttt{covered} region. The \texttt{missed} count ($\mathrm{condition} = \text{False} \land \mathrm{success} = \text{False}$) is dominated by early trials ($L = 1$) under high noise.

We extend the same protocol to larger semiprimes $N \in \{1147, 2491, 4087\}$ with a representative 4-noise subset and 100 trials per setup, yielding $1{,}200$ further measurements with zero violations (Appendix~\ref{app:large}). Per-noise-model breakdown for $N = 77$ is in Appendix~\ref{app:tables}. See \texttt{demo.py -{}-verify <N>} and \texttt{verify\_large\_run.py} in the companion code to reproduce.

\section{Related Work}
\label{sec:related}

We outline how Theorem~\ref{thm:c-det} relates to existing literature, including which ingredients are folklore and where the gap (if any) lies.

\paragraph{Knill (1995)~\cite{knill1995}.}
Proposes running the quantum step twice for the same base and taking the lcm of denominators recovered from convergents. This is \emph{single-base} lcm. Our $(C)$ extends this to multi-base lcm and augments the candidate pool with $\mathrm{divisors}(L)$.

\paragraph{McAnally (2001)~\cite{mcanally2001}.}
Same-base accumulation: enlarges the Fourier transform modulus to $Q \approx 2 w n^3$ and enumerates \emph{all} convergent denominators in a single measurement, achieving near-certainty per run. Same-base structure, different mechanism ($Q$ enlargement vs.\ $L$ accumulation across bases).

\paragraph{Bourdon--Williams (2007)~\cite{bourdonwilliams2007}.}
Sharp asymptotic lower bound on per-measurement success probability ($\sim 94\%$) via Carmichael-based analysis. Theoretical bound; no algorithmic change.

\paragraph{Ekerå (2021)~\cite{ekera2021}.}
Shows the order of a \emph{single} base typically suffices for complete factorization via smoothness extension. Different direction from our multi-base accumulation; complementary.

\paragraph{Quantum Carmichael (2021)~\cite{carmichaelquantum2021}.}
Algorithm~1 of this paper essentially runs the multi-base lcm to recover $\lambda(N)$. Closest prior work to our framework. The paper \emph{does not augment subsequent measurements' post-processing with $\mathrm{divisors}(L)$}; each order is computed independently.

\paragraph{Ekerå (2024)~\cite{ekera2024}.}
Comprehensive survey of post-processing techniques: continued fractions, lattice post-processing, offset search, divisor search of recovered $r/d$, Seifert's joint solving. \emph{Multi-base accumulation across different $a$ values is not surveyed as a post-processing technique here.}

\paragraph{Bach--Shallit (1996)~\cite{bachshallit1996}.}
Foundational textbook coverage of $r \mid \lambda(N)$, multi-element lcm convergence to $\lambda(N)$, and factoring from a known exponent (Miller--Rabin reduction). Related classical analysis of the Carmichael function distribution: Erd\H{o}s--Pomerance--Schmutz~\cite{erdosps1991}.

\paragraph{Pomerance et al.\ (2017)~\cite{pomerance2017}.}
Shows $e(G) \leq d + 2.752\ldots$ where $d$ is the maximum Sylow generator count. For semiprime $(\mathbb{Z}/N\mathbb{Z})^*$, this gives constant-bounded expected number of bases to \emph{generate} the group (hence reach $\lambda(N)$ via lcm).

\paragraph{Position of our contribution.}
Every individual ingredient --- $r \mid \lambda(N)$, multi-base lcm, divisor search given exponent, continued fractions, factoring from exponent --- is folklore or established in the references above. The \emph{clean separation theorem} (noise invariance via the classical verification gate, formalized as Theorem~\ref{thm:c-det}) and its corollaries appear to not be stated explicitly. We propose Theorem~\ref{thm:c-det} as the \emph{folklore packaged} result that gives a uniform explanation for the empirical robustness of multi-base post-processing.

\section{Limitations and Discussion}
\label{sec:limitations}

\paragraph{No measurement saving from adaptive base selection.}
A natural attempt --- pre-filter bases by $a^L \not\equiv 1 \bmod N$ so that every measured base extends $L$ --- yields no improvement in measurement count over the random-selection $+$ fast-path scheme (\texttt{shor\_quantum\_adaptive} vs.\ \texttt{shor\_quantum\_multi} in the companion code). The fast path already skips classical-only bases; pre-filtering shifts the work but does not eliminate any measurement.

\paragraph{Marginal value of lattice post-processing.}
The Knill--Mosca lattice approach strengthens single-base, multi-measurement convergent recovery. In our multi-base framework, $\mathrm{lcm}(r_{a_1}, \ldots, r_{a_k})$ already accomplishes equivalent recovery across bases. We did not find a regime where additional lattice basis reduction (LLL or simpler) provides measurable benefit over the lcm strategy.

\paragraph{The theorem is conditional on $r_a \mid L$.}
When $L = 1$ (start of algorithm), the condition fails and $(C) \equiv (B)$. The transient before $L$ reaches $\lambda(N)$ is governed by classical number theory: for typical semiprime $N = pq$, $\mathbb{E}[K_\lambda] = O(\log \log N)$ bases suffice (essentially proven by Pomerance et al.~\cite{pomerance2017} and the Carmichael-function paper~\cite{carmichaelquantum2021}).

\paragraph{Classical verification gate is assumed exact.}
Theorem~\ref{thm:c-det} relies on $\texttt{pow}(a, d, N)$ returning the correct value. On real hardware, this could in principle be corrupted by classical computation errors, but at standard machine precision this is negligible compared to quantum noise.

\paragraph{Scale of verification.}
Our experiments span $N \in \{77, 143, 209\}$ with 500 trials per noise setup (Sec.~\ref{sec:empirical}) and $N \in \{1147, 2491, 4087\}$ with 100 trials per setup (Appendix~\ref{app:large}), totaling $17{,}700$ measurements. The largest $N = 4087 = 61 \cdot 67$ requires a 24-qubit counting register ($Q = 2^{24} \approx 16.8\text{M}$ amplitudes, $256$ MB at \texttt{complex128}) --- at the upper end of single-machine state-vector simulation. Scaling to cryptographic $N$ requires no algorithmic change; the remaining bottleneck is pure classical simulation memory. Hardware verification on existing NISQ devices is a natural next step.

\section{Code and Reproducibility}
\label{sec:code}

All experiments are reproducible via the companion repository at\\
\url{https://github.com/Hashevolution/shor}.

\begin{lstlisting}[style=pseudostyle]
pip install numpy
git clone https://github.com/Hashevolution/shor
cd shor
python demo.py --verify 77 143 209          # reproduces Sec. 4 table
python verify_large_run.py 1147 2491 4087   # reproduces Appendix C (large N)
python demo.py --noise3 77 143              # 3-noise comparison
python demo.py --adaptive 77 143 209 323    # negative result, Sec. 6
\end{lstlisting}

\noindent
The codebase is approximately 700 lines of \texttt{numpy}. No quantum libraries are required.

\section{Conclusion}
\label{sec:conclusion}

We formalize a folklore observation about classical post-processing of multi-base Shor order-finding measurements: augmenting continued-fraction recovery with a divisor search over the accumulated lcm $L$ makes the procedure deterministic whenever the order of the current base divides $L$ --- independent of the measurement distribution, hence robust to all measurement-layer noise models. We verify this across $17{,}700$ trials at six composite sizes (zero violations), and show that two natural attempts to extend the result (adaptive base selection, lattice post-processing) yield no further measurement savings, indicating Theorem~\ref{thm:c-det} captures the essential structure at this scale.

The work is best viewed as expository and a foundation for explorations of more substantive improvements: noise-resilient implementations on hardware, the \texttt{lucky} region under structured measurement distributions, and connections to lattice-based factoring~\cite{regev2023}.

\begin{thebibliography}{12}

\bibitem{shor1994}
P.~W. Shor.
\newblock Algorithms for quantum computation: discrete logarithms and factoring.
\newblock In \emph{Proceedings 35th Annual Symposium on Foundations of Computer Science}, 1994.

\bibitem{knill1995}
E.~Knill.
\newblock On Shor's quantum factor finding algorithm: Increasing the probability of success and tradeoffs involving the {F}ourier transform modulus.
\newblock Technical report, Los Alamos National Laboratory, 1995.

\bibitem{mcanally2001}
D.~McAnally.
\newblock A refinement of {S}hor's algorithm.
\newblock arXiv:quant-ph/0112055, 2001.

\bibitem{nielsenchuang2000}
M.~A. Nielsen and I.~L. Chuang.
\newblock \emph{Quantum Computation and Quantum Information}.
\newblock Cambridge University Press, 2000.

\bibitem{bourdonwilliams2007}
P.~S. Bourdon and H.~T. Williams.
\newblock Sharp probability estimates for {S}hor's order-finding algorithm.
\newblock arXiv:quant-ph/0607148, 2007.

\bibitem{ekera2021}
M.~Ekerå.
\newblock On completely factoring any integer efficiently in a single run of an order-finding algorithm.
\newblock \emph{Quantum Information Processing}, 20(6), 2021.
\newblock arXiv:2007.10044.

\bibitem{ekera2024}
M.~Ekerå.
\newblock On the success probability of quantum order finding.
\newblock \emph{ACM Transactions on Quantum Computing}, 2024.
\newblock arXiv:2201.07791.

\bibitem{bachshallit1996}
E.~Bach and J.~Shallit.
\newblock \emph{Algorithmic Number Theory}, volume~1.
\newblock MIT Press, 1996.

\bibitem{erdosps1991}
P.~Erdős, C.~Pomerance, and E.~Schmutz.
\newblock {C}armichael's lambda function.
\newblock \emph{Acta Arithmetica}, 58(4):363--385, 1991.

\bibitem{pomerance2017}
C.~Pomerance et~al.
\newblock The expected number of elements to generate a finite group with $d$-generated {S}ylow subgroups.
\newblock arXiv:1707.07193, 2017.

\bibitem{regev2023}
O.~Regev.
\newblock An efficient quantum factoring algorithm.
\newblock arXiv:2308.06572, 2023.

\bibitem{carmichaelquantum2021}
A quantum algorithm for computing the {C}armichael function.
\newblock arXiv:2111.02488, 2021.

\end{thebibliography}

\appendix

\section{Reference Implementation}
\label{app:impl}

\begin{lstlisting}[style=pythonstyle]
def divisors(n):
    """Divisors of n, ascending."""
    small, large = [], []
    i = 1
    while i*i <= n:
        if n % i == 0:
            small.append(i)
            if i != n // i:
                large.append(n // i)
        i += 1
    return small + large[::-1]


def convergent_denominators(k, Q, max_denom):
    """All convergent denominators of k/Q up to max_denom."""
    a, b = k, Q
    h_prev, h_curr = 0, 1
    k_prev, k_curr = 1, 0
    denoms = []
    while b != 0:
        q = a // b
        a, b = b, a - q*b
        h_prev, h_curr = h_curr, q*h_curr + h_prev
        k_prev, k_curr = k_curr, q*k_curr + k_prev
        if 0 < k_curr <= max_denom:
            denoms.append(k_curr)
        elif k_curr > max_denom:
            break
    return denoms


def minimize_order(a, N, candidate):
    if candidate <= 0 or pow(a, candidate, N) != 1:
        return 0
    r = candidate
    for p in prime_factors(r):
        while r % p == 0 and pow(a, r // p, N) == 1:
            r //= p
    return r


def C(a, N, k, Q, L):
    """The (C) post-processing procedure."""
    cands = set(convergent_denominators(k, Q, N - 1))
    if L > 1:
        cands.update(divisors(L))
    valid = [d for d in cands if d > 0 and pow(a, d, N) == 1]
    if not valid:
        return 0
    return minimize_order(a, N, min(valid))
\end{lstlisting}

\section{Per-Noise-Model Violation Counts ($N = 77$)}
\label{app:tables}

\begin{center}
\begin{tabular}{lrrrr}
\toprule
Noise & covered/500 & violations & lucky & missed \\
\midrule
noise-free               & 493 & 0 & 3 & 4 \\
depolarizing $p=0.3$     & 499 & 0 & 1 & 0 \\
depolarizing $p=0.8$     & 494 & 0 & 3 & 3 \\
readout flip $p=0.3$     & 487 & 0 & 3 & 10 \\
bias zero $p=0.5$        & 499 & 0 & 1 & 0 \\
phase $\sigma=1.0$       & 496 & 0 & 3 & 1 \\
phase $\sigma=2.5$       & 481 & 0 & 2 & 17 \\
amp.\ damp.\ $\gamma=0.01$ & 491 & 0 & 2 & 7 \\
amp.\ damp.\ $\gamma=0.05$ & 492 & 0 & 3 & 5 \\
modexp $q=0.3$           & 493 & 0 & 3 & 4 \\
modexp $q=0.8$           & 476 & 0 & 3 & 21 \\
\bottomrule
\end{tabular}
\end{center}

\noindent
Identical pattern (violations $= 0$) for $N = 143$ and $N = 209$; see \texttt{demo.py -{}-verify} output in the companion code.

\section{Extended Verification at Larger $N$}
\label{app:large}

We additionally verify Theorem~\ref{thm:c-det} at three larger semiprimes, with a representative subset of noise models (one each from categories A, B, C plus noise-free) and reduced trial count (100 per setup) due to the increased per-measurement cost ($\approx 0.45$\,s at $N = 1147$, $\approx 1.9$\,s at $N = 4087$ on a single CPU).

\begin{center}
\begin{tabular}{lrrrr}
\toprule
$N$ / Noise & covered/100 & violations & lucky & missed \\
\midrule
\multicolumn{5}{l}{\emph{$N = 1147 = 31 \cdot 37$}} \\
\quad noise-free          & 91 & 0 & 2 &  7 \\
\quad depol $p=0.8$       & 93 & 0 & 3 &  4 \\
\quad phase $\sigma=2.5$  &  9 & 0 & 4 & 87 \\
\quad modexp $q=0.8$      & 61 & 0 & 2 & 37 \\
\midrule
\multicolumn{5}{l}{\emph{$N = 2491 = 47 \cdot 53$}} \\
\quad noise-free          & 97 & 0 & 2 &  1 \\
\quad depol $p=0.8$       & 96 & 0 & 1 &  3 \\
\quad phase $\sigma=2.5$  &  1 & 0 & 1 & 98 \\
\quad modexp $q=0.8$      & 59 & 0 & 2 & 39 \\
\midrule
\multicolumn{5}{l}{\emph{$N = 4087 = 61 \cdot 67$}} \\
\quad noise-free          & 88 & 0 & 3 &  9 \\
\quad depol $p=0.8$       & 93 & 0 & 2 &  5 \\
\quad phase $\sigma=2.5$  & 50 & 0 & 3 & 47 \\
\quad modexp $q=0.8$      &  1 & 0 & 1 & 98 \\
\bottomrule
\end{tabular}
\end{center}

\noindent
\textbf{Total: $1{,}200$ additional measurements, $0$ violations.} Combined with Sec.~\ref{sec:empirical}, the verification covers $17{,}700$ measurements.

Under extreme noise the \texttt{covered} fraction can drop sharply --- \texttt{phase $\sigma = 2.5$} at $N = 2491$ and \texttt{modexp $q = 0.8$} at $N = 4087$ both yield $\texttt{covered} = 1$ --- because the measurement is so degraded that $L$ cannot accumulate. Theorem~\ref{thm:c-det} then applies only vacuously on most trials, but on those where it does apply, success remains $100\%$. The non-monotonicity across $N$ for the same noise (\texttt{phase $\sigma = 2.5$} covered: $9, 1, 50$ at $N = 1147, 2491, 4087$) is a sample-size artifact at 100 trials and reflects which random seeds happened to land on bases with small $r_a$ early.

Reproduce: \texttt{python verify\_large\_run.py 1147 2491 4087} ($\approx 26$ min total on a single CPU).

\end{document}
